\begin{table}[ht]
\centering
\footnotesize
\caption{主对比 P1--P3 的效应量与推断结果。置信区间为边际 95\% basic 自助区间；原始单侧 bootstrap $p$ 值由 10,000 次整簇重采样的尾部计数按加一规则计算，随后对三项主对比实施 Holm 逐步校正，表中报告调整后 $p$ 值。}\label{tab:primary}
\begin{tabular}{llccccc}
\toprule
对比 & 估计量 & 点估计 & 阈值 & 超额效应 & 95\% basic CI & Holm $p$\\
\midrule
P1 & fixed-$M$：$s{=}16$ 减 $s{=}64$ & 0.5354 & 0.10 & 0.4354 & $[0.5210,0.5494]$ & 0.00030\\
P2 & fixed-$T$：$s{=}16$ 减 $s{=}64$ & 0.5330 & 0.10 & 0.4330 & $[0.5194,0.5462]$ & 0.00030\\
P3 & $0.05$ 减 M1 & 0.0266 & 0.01 & 0.0166 & $[0.0224,0.0308]$ & 0.00030\\
\bottomrule
\end{tabular}

\begin{minipage}{0.94\textwidth}
\vspace{0.4em}
\footnotesize 注：P1/P2 为 M0 拒绝率差（P1 固定窗口数，P2 固定观测长度），P3 为 $0.05$ 减 M1 拒绝率。三项尾部计数经真实 bootstrap 抽样核验均为 0，故原始 $p=1/10001\approx9.999\times10^{-5}$；经 Holm 逐步校正后三项调整 $p$ 均为 $3/10001\approx0.00030$。点估计与阈值之差为超额效应；正文按预设单侧阈值与多重性规则给出支持与否的表述，表中仅列统计量。
\end{minipage}
\end{table}
